% --- Archon physics formalization source begin ---
% archon:physics
% archon:covers PhyXMiniProblems/problem_phyx_mini_0811.lean
% archon:source-report reports/phyx_mini/problem_phyx_mini_0811.source.json

\chapter{Physics problem phyx\_mini\_0811}
\label{ch:PhyXMiniProblems_problem_phyx_mini_0811}

\paragraph{Problem source.}
\#\# Physical scenario

Your cousin Throckmorton skateboards from rest down a curved, frictionless ramp. If we treat Throcky and his skateboard as a particle, he moves through a quarter-circle with radius \textbackslash{}( R = 3.00 \textbackslash{}) m as shown in figure. Throcky and his skateboard have a total mass of 25.0 kg.

\#\# Question

Find his speed at the bottom of the ramp.

\#\# Answer choices

- A: A: 5.42m/s

- B: B: 7.67m/s

- C: C: 4.98m/s

- D: D: 9.02m/s

\#\# Figure caption (auxiliary; use the image as primary evidence)

The image depicts a scenario of a skateboarder on a curved ramp, illustrating principles of physics. Here’s a breakdown of the elements:

1. **Ramp**: A smooth, curved surface is shown with a concave profile, typically representing a half-pipe used in skateboarding.

2. **Skateboarder**: The person is illustrated in three different positions along the ramp, indicating motion. 

3. **Positions**:
   - **Point 1**: At the top-left of the ramp. The skater is at rest (\textbackslash{}(v\_1 = 0\textbackslash{})) with \textbackslash{}(y\_1 = R\textbackslash{}). 
   - **Point 2**: At the bottom of the ramp. The skater is standing with speed \textbackslash{}(v\_2\textbackslash{}), and \textbackslash{}(y\_2 = 0\textbackslash{}).

4. **Text and Labels**:
   - **Point 1**: Annotated with \textbackslash{}(y\_1 = R\textbackslash{}) and \textbackslash{}(v\_1 = 0\textbackslash{}). 
   - **Point 2**: Annotated with \textbackslash{}(y\_2 = 0\textbackslash{}) and \textbackslash{}(v\_2\textbackslash{}).
   - **O**: The center of curvature of the ramp, marked with a black dot.
   - **R = 3.00 m**: Indicates the radius of curvature of the ramp.

5. **Arrows**:
   - A dashed blue arrow traces the path of the skateboarder from Point 1 to Point 2.
   - A solid green arrow at Point 2 indicates the direction of the skateboarder's velocity (\textbackslash{}(v\_2\textbackslash{})) at the bottom.

This image likely illustrates concepts of energy conservation and motion dynamics on a frictionless surface.

\#\# Dataset metadata

Work and Energy; Physical Model Grounding Reasoning, Multi-Formula Reasoning

\paragraph{Recorded answer/context.}
B: B: 7.67m/s

\paragraph{Figure/image path.}
/root/proposal\_for\_physic/science-mango/phyx\_mini\_run/phyx\_data/test\_image/811.png

\paragraph{Formalization target.}
create a compiling Lean file with sorry bodies at `PhyXMiniProblems/problem\_phyx\_mini\_0811.lean`. The Lean declarations must preserve the physical quantities, dimensions or dimensional roles, figure labels, governing-law hypotheses, and final relation expressed by this problem.
Use Mathlib/Physlib names found through LeanExplore where available. If a physics API is missing, introduce faithful local abstractions rather than scalar placeholder aliases.

\begin{theorem}[Physics formalization target]
\label{thm:physics:phyx_mini_0811:target}
\lean{PhyXMiniProblems.ProblemPhyXMini0811.skateboarderBottomSpeed}
\uses{def:physics:phyx-mini-0811:phyxminiproblems-problemphyxmini0811-speedinmeterspersecond, def:physics:phyx-mini-0811:phyxminiproblems-problemphyxmini0811-skateboardrampsetup, def:physics:phyx-mini-0811:phyxminiproblems-problemphyxmini0811-matchesskateboardrampscenario, def:physics:phyx-mini-0811:phyxminiproblems-problemphyxmini0811-matchesgivenproblemdata, def:physics:phyx-mini-0811:phyxminiproblems-problemphyxmini0811-matchesprimaryskateboardrampfigure, def:physics:phyx-mini-0811:phyxminiproblems-problemphyxmini0811-hasphysicalskateboardrampparameters, def:physics:phyx-mini-0811:phyxminiproblems-problemphyxmini0811-satisfiesskateboardrampenergylaws, def:physics:phyx-mini-0811:phyxminiproblems-problemphyxmini0811-agreeswithhundredthspeedreadout, def:physics:phyx-mini-0811:phyxminiproblems-problemphyxmini0811-answerchoicereportsbottomspeed}
The assigned autoformalize agent should translate this physics problem into Lean declarations in the covered file, with theorem and lemma proof bodies written as `by sorry`.
\end{theorem}
\begin{proof}
This is an autoformalization task, not a proof task. Produce statements that can later be proved by the physics prover without weakening the physical model.
\end{proof}

% --- Archon declaration topology begin ---
\section{Declaration topology}
\begin{definition}[acceleration Dimension]
  \label{def:physics:phyx-mini-0811:phyxminiproblems-problemphyxmini0811-accelerationdimension}
  \lean{PhyXMiniProblems.ProblemPhyXMini0811.accelerationDimension}
  Linear acceleration has physical dimension L T⁻²
\end{definition}
\begin{proof}
  This is the declared model definition of acceleration Dimension.
\end{proof}
\begin{definition}[Mass Quantity]
  \label{def:physics:phyx-mini-0811:phyxminiproblems-problemphyxmini0811-massquantity}
  \lean{PhyXMiniProblems.ProblemPhyXMini0811.MassQuantity}
  A nonnegative, unit-independent physical mass
\end{definition}
\begin{proof}
  This is the declared model definition of Mass Quantity.
\end{proof}
\begin{definition}[Length Quantity]
  \label{def:physics:phyx-mini-0811:phyxminiproblems-problemphyxmini0811-lengthquantity}
  \lean{PhyXMiniProblems.ProblemPhyXMini0811.LengthQuantity}
  A nonnegative, unit-independent physical length or height
\end{definition}
\begin{proof}
  This is the declared model definition of Length Quantity.
\end{proof}
\begin{definition}[Speed Quantity]
  \label{def:physics:phyx-mini-0811:phyxminiproblems-problemphyxmini0811-speedquantity}
  \lean{PhyXMiniProblems.ProblemPhyXMini0811.SpeedQuantity}
  A nonnegative speed magnitude
\end{definition}
\begin{proof}
  This is the declared model definition of Speed Quantity.
\end{proof}
\begin{definition}[Acceleration Quantity]
  \label{def:physics:phyx-mini-0811:phyxminiproblems-problemphyxmini0811-accelerationquantity}
  \lean{PhyXMiniProblems.ProblemPhyXMini0811.AccelerationQuantity}
  \uses{def:physics:phyx-mini-0811:phyxminiproblems-problemphyxmini0811-accelerationdimension}
  A nonnegative gravitational-acceleration magnitude
\end{definition}
\begin{proof}
  This is the declared model definition of Acceleration Quantity.
\end{proof}
\begin{definition}[Energy Quantity]
  \label{def:physics:phyx-mini-0811:phyxminiproblems-problemphyxmini0811-energyquantity}
  \lean{PhyXMiniProblems.ProblemPhyXMini0811.EnergyQuantity}
  A signed work or energy quantity
\end{definition}
\begin{proof}
  This is the declared model definition of Energy Quantity.
\end{proof}
\begin{definition}[mass In Kilograms]
  \label{def:physics:phyx-mini-0811:phyxminiproblems-problemphyxmini0811-massinkilograms}
  \lean{PhyXMiniProblems.ProblemPhyXMini0811.massInKilograms}
  \uses{def:physics:phyx-mini-0811:phyxminiproblems-problemphyxmini0811-massquantity}
  Read a physical mass in coherent-SI kilograms
\end{definition}
\begin{proof}
  This is the declared model definition of mass In Kilograms.
\end{proof}
\begin{definition}[length In Meters]
  \label{def:physics:phyx-mini-0811:phyxminiproblems-problemphyxmini0811-lengthinmeters}
  \lean{PhyXMiniProblems.ProblemPhyXMini0811.lengthInMeters}
  \uses{def:physics:phyx-mini-0811:phyxminiproblems-problemphyxmini0811-lengthquantity}
  Read a physical length in coherent-SI metres
\end{definition}
\begin{proof}
  This is the declared model definition of length In Meters.
\end{proof}
\begin{definition}[speed In Meters Per Second]
  \label{def:physics:phyx-mini-0811:phyxminiproblems-problemphyxmini0811-speedinmeterspersecond}
  \lean{PhyXMiniProblems.ProblemPhyXMini0811.speedInMetersPerSecond}
  \uses{def:physics:phyx-mini-0811:phyxminiproblems-problemphyxmini0811-speedquantity}
  Read a speed magnitude in coherent-SI metres per second
\end{definition}
\begin{proof}
  This is the declared model definition of speed In Meters Per Second.
\end{proof}
\begin{definition}[acceleration In Meters Per Second Squared]
  \label{def:physics:phyx-mini-0811:phyxminiproblems-problemphyxmini0811-accelerationinmeterspersecondsquared}
  \lean{PhyXMiniProblems.ProblemPhyXMini0811.accelerationInMetersPerSecondSquared}
  \uses{def:physics:phyx-mini-0811:phyxminiproblems-problemphyxmini0811-accelerationquantity}
  Read an acceleration magnitude in coherent-SI metres per second squared
\end{definition}
\begin{proof}
  This is the declared model definition of acceleration In Meters Per Second Squared.
\end{proof}
\begin{definition}[energy In Joules]
  \label{def:physics:phyx-mini-0811:phyxminiproblems-problemphyxmini0811-energyinjoules}
  \lean{PhyXMiniProblems.ProblemPhyXMini0811.energyInJoules}
  \uses{def:physics:phyx-mini-0811:phyxminiproblems-problemphyxmini0811-energyquantity}
  Read work or energy in coherent-SI joules
\end{definition}
\begin{proof}
  This is the declared model definition of energy In Joules.
\end{proof}
\begin{definition}[Ramp Point]
  \label{def:physics:phyx-mini-0811:phyxminiproblems-problemphyxmini0811-ramppoint}
  \lean{PhyXMiniProblems.ProblemPhyXMini0811.RampPoint}
  The two labeled endpoints used in the energy comparison
\end{definition}
\begin{proof}
  This is the declared model definition of Ramp Point.
\end{proof}
\begin{definition}[Body Model]
  \label{def:physics:phyx-mini-0811:phyxminiproblems-problemphyxmini0811-bodymodel}
  \lean{PhyXMiniProblems.ProblemPhyXMini0811.BodyModel}
  The body idealization stated in the problem
\end{definition}
\begin{proof}
  This is the declared model definition of Body Model.
\end{proof}
\begin{definition}[Ramp Geometry]
  \label{def:physics:phyx-mini-0811:phyxminiproblems-problemphyxmini0811-rampgeometry}
  \lean{PhyXMiniProblems.ProblemPhyXMini0811.RampGeometry}
  The geometric class of the pictured ramp segment
\end{definition}
\begin{proof}
  This is the declared model definition of Ramp Geometry.
\end{proof}
\begin{definition}[Surface Interaction]
  \label{def:physics:phyx-mini-0811:phyxminiproblems-problemphyxmini0811-surfaceinteraction}
  \lean{PhyXMiniProblems.ProblemPhyXMini0811.SurfaceInteraction}
  Tangential interaction between the skateboard and the ramp
\end{definition}
\begin{proof}
  This is the declared model definition of Surface Interaction.
\end{proof}
\begin{definition}[Horizontal Direction]
  \label{def:physics:phyx-mini-0811:phyxminiproblems-problemphyxmini0811-horizontaldirection}
  \lean{PhyXMiniProblems.ProblemPhyXMini0811.HorizontalDirection}
  Horizontal directions used by the bottom velocity arrow
\end{definition}
\begin{proof}
  This is the declared model definition of Horizontal Direction.
\end{proof}
\begin{definition}[Figure Object]
  \label{def:physics:phyx-mini-0811:phyxminiproblems-problemphyxmini0811-figureobject}
  \lean{PhyXMiniProblems.ProblemPhyXMini0811.FigureObject}
  Physical or graphical objects visibly present in image 811.png
\end{definition}
\begin{proof}
  This is the declared model definition of Figure Object.
\end{proof}
\begin{definition}[Figure Label]
  \label{def:physics:phyx-mini-0811:phyxminiproblems-problemphyxmini0811-figurelabel}
  \lean{PhyXMiniProblems.ProblemPhyXMini0811.FigureLabel}
  Literal semantic labels printed in the primary image
\end{definition}
\begin{proof}
  This is the declared model definition of Figure Label.
\end{proof}
\begin{definition}[Skateboard Ramp Figure]
  \label{def:physics:phyx-mini-0811:phyxminiproblems-problemphyxmini0811-skateboardrampfigure}
  \lean{PhyXMiniProblems.ProblemPhyXMini0811.SkateboardRampFigure}
  \uses{def:physics:phyx-mini-0811:phyxminiproblems-problemphyxmini0811-ramppoint, def:physics:phyx-mini-0811:phyxminiproblems-problemphyxmini0811-horizontaldirection, def:physics:phyx-mini-0811:phyxminiproblems-problemphyxmini0811-figureobject, def:physics:phyx-mini-0811:phyxminiproblems-problemphyxmini0811-figurelabel}
  Literal evidence transcribed from the supplied bitmap. Its real fields are printed labels, not definitions of the physical trajectory or unknown speed
\end{definition}
\begin{proof}
  This is the declared model definition of Skateboard Ramp Figure.
\end{proof}
\begin{definition}[Skateboard Ramp Setup]
  \label{def:physics:phyx-mini-0811:phyxminiproblems-problemphyxmini0811-skateboardrampsetup}
  \lean{PhyXMiniProblems.ProblemPhyXMini0811.SkateboardRampSetup}
  \uses{def:physics:phyx-mini-0811:phyxminiproblems-problemphyxmini0811-massquantity, def:physics:phyx-mini-0811:phyxminiproblems-problemphyxmini0811-lengthquantity, def:physics:phyx-mini-0811:phyxminiproblems-problemphyxmini0811-speedquantity, def:physics:phyx-mini-0811:phyxminiproblems-problemphyxmini0811-accelerationquantity, def:physics:phyx-mini-0811:phyxminiproblems-problemphyxmini0811-energyquantity, def:physics:phyx-mini-0811:phyxminiproblems-problemphyxmini0811-ramppoint, def:physics:phyx-mini-0811:phyxminiproblems-problemphyxmini0811-bodymodel, def:physics:phyx-mini-0811:phyxminiproblems-problemphyxmini0811-rampgeometry, def:physics:phyx-mini-0811:phyxminiproblems-problemphyxmini0811-surfaceinteraction, def:physics:phyx-mini-0811:phyxminiproblems-problemphyxmini0811-skateboardrampfigure}
  The bottom speed is an independent observable field. In particular, it is not defined from an answer choice or from the desired energy formula
\end{definition}
\begin{proof}
  This is the declared model definition of Skateboard Ramp Setup.
\end{proof}
\begin{definition}[Matches Skateboard Ramp Scenario]
  \label{def:physics:phyx-mini-0811:phyxminiproblems-problemphyxmini0811-matchesskateboardrampscenario}
  \lean{PhyXMiniProblems.ProblemPhyXMini0811.MatchesSkateboardRampScenario}
  \uses{def:physics:phyx-mini-0811:phyxminiproblems-problemphyxmini0811-speedinmeterspersecond, def:physics:phyx-mini-0811:phyxminiproblems-problemphyxmini0811-skateboardrampsetup}
  The prose model: a particle descends a frictionless quarter-circle
\end{definition}
\begin{proof}
  This is the declared model definition of Matches Skateboard Ramp Scenario.
\end{proof}
\begin{definition}[Matches Given Problem Data]
  \label{def:physics:phyx-mini-0811:phyxminiproblems-problemphyxmini0811-matchesgivenproblemdata}
  \lean{PhyXMiniProblems.ProblemPhyXMini0811.MatchesGivenProblemData}
  \uses{def:physics:phyx-mini-0811:phyxminiproblems-problemphyxmini0811-massinkilograms, def:physics:phyx-mini-0811:phyxminiproblems-problemphyxmini0811-lengthinmeters, def:physics:phyx-mini-0811:phyxminiproblems-problemphyxmini0811-accelerationinmeterspersecondsquared, def:physics:phyx-mini-0811:phyxminiproblems-problemphyxmini0811-skateboardrampsetup}
  Numerical problem data and the standard near-Earth gravitational calibration. The gravity value is a model input, not a conclusion about the bottom speed
\end{definition}
\begin{proof}
  This is the declared model definition of Matches Given Problem Data.
\end{proof}
\begin{definition}[Matches Primary Skateboard Ramp Figure]
  \label{def:physics:phyx-mini-0811:phyxminiproblems-problemphyxmini0811-matchesprimaryskateboardrampfigure}
  \lean{PhyXMiniProblems.ProblemPhyXMini0811.MatchesPrimarySkateboardRampFigure}
  \uses{def:physics:phyx-mini-0811:phyxminiproblems-problemphyxmini0811-lengthinmeters, def:physics:phyx-mini-0811:phyxminiproblems-problemphyxmini0811-speedinmeterspersecond, def:physics:phyx-mini-0811:phyxminiproblems-problemphyxmini0811-skateboardrampsetup}
  Facts read from the primary image: the dashed path joins Point 1 to Point 2, the radius is drawn from O, and the unknown v₂ arrow is tangent and points right. The endpoint labels are connected to the independent physical fields, but no numerical value is assigned to the Point 2 speed
\end{definition}
\begin{proof}
  This is the declared model definition of Matches Primary Skateboard Ramp Figure.
\end{proof}
\begin{definition}[Has Physical Skateboard Ramp Parameters]
  \label{def:physics:phyx-mini-0811:phyxminiproblems-problemphyxmini0811-hasphysicalskateboardrampparameters}
  \lean{PhyXMiniProblems.ProblemPhyXMini0811.HasPhysicalSkateboardRampParameters}
  \uses{def:physics:phyx-mini-0811:phyxminiproblems-problemphyxmini0811-massinkilograms, def:physics:phyx-mini-0811:phyxminiproblems-problemphyxmini0811-lengthinmeters, def:physics:phyx-mini-0811:phyxminiproblems-problemphyxmini0811-accelerationinmeterspersecondsquared, def:physics:phyx-mini-0811:phyxminiproblems-problemphyxmini0811-energyinjoules, def:physics:phyx-mini-0811:phyxminiproblems-problemphyxmini0811-skateboardrampsetup}
  Positivity and nondegeneracy conditions selecting the physical branch
\end{definition}
\begin{proof}
  This is the declared model definition of Has Physical Skateboard Ramp Parameters.
\end{proof}
\begin{definition}[total Mechanical Energy In Joules]
  \label{def:physics:phyx-mini-0811:phyxminiproblems-problemphyxmini0811-totalmechanicalenergyinjoules}
  \lean{PhyXMiniProblems.ProblemPhyXMini0811.totalMechanicalEnergyInJoules}
  \uses{def:physics:phyx-mini-0811:phyxminiproblems-problemphyxmini0811-energyinjoules, def:physics:phyx-mini-0811:phyxminiproblems-problemphyxmini0811-ramppoint, def:physics:phyx-mini-0811:phyxminiproblems-problemphyxmini0811-skateboardrampsetup}
  Total translational mechanical energy at a labeled endpoint
\end{definition}
\begin{proof}
  This is the declared model definition of total Mechanical Energy In Joules.
\end{proof}
\begin{definition}[Satisfies Skateboard Ramp Energy Laws]
  \label{def:physics:phyx-mini-0811:phyxminiproblems-problemphyxmini0811-satisfiesskateboardrampenergylaws}
  \lean{PhyXMiniProblems.ProblemPhyXMini0811.SatisfiesSkateboardRampEnergyLaws}
  \uses{def:physics:phyx-mini-0811:phyxminiproblems-problemphyxmini0811-massinkilograms, def:physics:phyx-mini-0811:phyxminiproblems-problemphyxmini0811-lengthinmeters, def:physics:phyx-mini-0811:phyxminiproblems-problemphyxmini0811-speedinmeterspersecond, def:physics:phyx-mini-0811:phyxminiproblems-problemphyxmini0811-accelerationinmeterspersecondsquared, def:physics:phyx-mini-0811:phyxminiproblems-problemphyxmini0811-energyinjoules, def:physics:phyx-mini-0811:phyxminiproblems-problemphyxmini0811-skateboardrampsetup, def:physics:phyx-mini-0811:phyxminiproblems-problemphyxmini0811-totalmechanicalenergyinjoules}
  The point-particle kinetic and gravitational potential energy formulae, together with endpoint work-energy balance. The contact-work term is retained explicitly so that its vanishing is tied to the frictionless ramp rather than hidden in the desired speed equation
\end{definition}
\begin{proof}
  This is the declared model definition of Satisfies Skateboard Ramp Energy Laws.
\end{proof}
\begin{definition}[Answer Choice]
  \label{def:physics:phyx-mini-0811:phyxminiproblems-problemphyxmini0811-answerchoice}
  \lean{PhyXMiniProblems.ProblemPhyXMini0811.AnswerChoice}
  The four speed choices displayed with the problem
\end{definition}
\begin{proof}
  This is the declared model definition of Answer Choice.
\end{proof}
\begin{definition}[displayed Speed In Meters Per Second]
  \label{def:physics:phyx-mini-0811:phyxminiproblems-problemphyxmini0811-displayedspeedinmeterspersecond}
  \lean{PhyXMiniProblems.ProblemPhyXMini0811.displayedSpeedInMetersPerSecond}
  \uses{def:physics:phyx-mini-0811:phyxminiproblems-problemphyxmini0811-answerchoice}
  Metres-per-second value printed beside each answer label
\end{definition}
\begin{proof}
  This is the declared model definition of displayed Speed In Meters Per Second.
\end{proof}
\begin{definition}[recorded Dataset Answer]
  \label{def:physics:phyx-mini-0811:phyxminiproblems-problemphyxmini0811-recordeddatasetanswer}
  \lean{PhyXMiniProblems.ProblemPhyXMini0811.recordedDatasetAnswer}
  \uses{def:physics:phyx-mini-0811:phyxminiproblems-problemphyxmini0811-answerchoice}
  The answer label stored in the source dataset; metadata, not a premise
\end{definition}
\begin{proof}
  This is the declared model definition of recorded Dataset Answer.
\end{proof}
\begin{definition}[Agrees With Hundredth Speed Readout]
  \label{def:physics:phyx-mini-0811:phyxminiproblems-problemphyxmini0811-agreeswithhundredthspeedreadout}
  \lean{PhyXMiniProblems.ProblemPhyXMini0811.AgreesWithHundredthSpeedReadout}
  \uses{def:physics:phyx-mini-0811:phyxminiproblems-problemphyxmini0811-speedquantity, def:physics:phyx-mini-0811:phyxminiproblems-problemphyxmini0811-speedinmeterspersecond}
  Agreement with a speed displayed to the nearest hundredth of a metre per second
\end{definition}
\begin{proof}
  This is the declared model definition of Agrees With Hundredth Speed Readout.
\end{proof}
\begin{definition}[Answer Choice Reports Bottom Speed]
  \label{def:physics:phyx-mini-0811:phyxminiproblems-problemphyxmini0811-answerchoicereportsbottomspeed}
  \lean{PhyXMiniProblems.ProblemPhyXMini0811.AnswerChoiceReportsBottomSpeed}
  \uses{def:physics:phyx-mini-0811:phyxminiproblems-problemphyxmini0811-skateboardrampsetup, def:physics:phyx-mini-0811:phyxminiproblems-problemphyxmini0811-answerchoice, def:physics:phyx-mini-0811:phyxminiproblems-problemphyxmini0811-displayedspeedinmeterspersecond, def:physics:phyx-mini-0811:phyxminiproblems-problemphyxmini0811-agreeswithhundredthspeedreadout}
  A displayed choice reports the independently modeled bottom speed
\end{definition}
\begin{proof}
  This is the declared model definition of Answer Choice Reports Bottom Speed.
\end{proof}
\begin{lemma}[bottom Speed Squared]
  \label{lem:physics:phyx-mini-0811:phyxminiproblems-problemphyxmini0811-bottomspeedsquared}
  \lean{PhyXMiniProblems.ProblemPhyXMini0811.bottomSpeedSquared}
  \uses{def:physics:phyx-mini-0811:phyxminiproblems-problemphyxmini0811-speedinmeterspersecond, def:physics:phyx-mini-0811:phyxminiproblems-problemphyxmini0811-skateboardrampsetup, def:physics:phyx-mini-0811:phyxminiproblems-problemphyxmini0811-matchesskateboardrampscenario, def:physics:phyx-mini-0811:phyxminiproblems-problemphyxmini0811-matchesgivenproblemdata, def:physics:phyx-mini-0811:phyxminiproblems-problemphyxmini0811-matchesprimaryskateboardrampfigure, def:physics:phyx-mini-0811:phyxminiproblems-problemphyxmini0811-hasphysicalskateboardrampparameters, def:physics:phyx-mini-0811:phyxminiproblems-problemphyxmini0811-satisfiesskateboardrampenergylaws}
  Energy conservation across the vertical drop R gives the exact squared bottom speed 2 g R = 294/5 (m/s)². This is a derived conclusion, not a law field or figure readout
\end{lemma}
\begin{proof}
  The conclusion follows from the governing relations and figure readouts named in the statement of bottom Speed Squared.
\end{proof}
% --- Archon declaration topology end ---
% --- Archon physics formalization source end ---
